\documentclass[10pt,a4paper]{amsart}
\usepackage[margin=19mm]{geometry}
\usepackage{amsmath,amssymb,mathtools}
\usepackage[scr=rsfs,cal=euler]{mathalfa}
\usepackage{xfrac}
\usepackage[unicode,hidelinks,bookmarks=false]{hyperref}
\newcommand{\ie}{i.e.\ }
\newcommand{\cf}{cf.\ }
\newcommand{\resp}{respectively\ }
\newcommand{\Subsection}{\subsection}
\providecommand{\nobreakdash}{\nobreak\hbox{-}}
\setlength{\emergencystretch}{2em}
\multlinegap0pt
\usepackage{bm}
\usepackage{bbm}
\usepackage{dsfont}
\usepackage{xspace}
\usepackage{cancel}
\usepackage{mathtools}
\usepackage{amsthm}
\makeatletter
\@ifundefined{theorem}{\newtheorem{theorem}{Theorem}}{}
\@ifundefined{claim}{\newtheorem{claim}[theorem]{Claim}}{}
\@ifundefined{fact}{\newtheorem{fact}[theorem]{Fact}}{}
\@ifundefined{proposition}{\newtheorem{proposition}[theorem]{Proposition}}{}
\makeatother
\newcommand{\forkindep}[1][]{%
  \mathrel{
    \mathop{
      \vcenter{
        \hbox{\oalign{\noalign{\kern-.3ex}\hfil$\vert$\hfil\cr
              \noalign{\kern-.7ex}
              $\smile$\cr\noalign{\kern-.3ex}}}
      }
    }\displaylimits_{#1}
  }
}
\newcommand{\vd}{\mathrm{OD}}
\title[Independence relations in the Solovay model II]{Independence relations in the Solovay model II}
\author{Jindrich Zapletal}
\date{Source: arXiv:2606.27011v1 (CC BY 4.0)}
\begin{document}
\maketitle

\noindent\emph{Source and licence.} Independence relations in the Solovay model II. Version arXiv:2606.27011v1 (CC BY 4.0), distributed under the Creative Commons Attribution 4.0 International licence, as recorded in the arXiv licence metadata for this exact version and shown on the abstract page https://arxiv.org/abs/2606.27011v1. Full rights notice: licenses/arxiv-2606.27011-CC-BY-4.0.txt.\par\smallskip

\noindent\textit{Editorial scope.} Counted unit: the Proposition whose source label is \texttt{ab} in the article (source0.tex lines 752--755, source label \texttt{ab}), delivered together with the article's own complete original proof of it (source0.tex lines 757--780, one \texttt{proof} environment of 24 lines). No mathematics is written, repaired or re-derived: statement and proof are the article's own text line for line. Required context retained uncounted: fact1 (lines 919--929). Every definition, display and label of those lines is retained unchanged. Omitted from this package and not redistributed: 1--751, 756--756, 781--918, 930--971. Nothing inside a retained excerpt is omitted or altered: every retained body is delivered exactly as the article's lines are written, so no omission receipt of any kind accompanies this package. Citations: the retained excerpts carry no citation command at all, so no bibliography entry of the article is transcribed and no reference is redistributed. Reference adaptation: none was needed and none was made; every reference inside the delivered unit resolves inside it, so no unresolved reference or citation is produced. Printed slips kept literal: no printed word, symbol, hypothesis or number of the counted statement or of its proof was altered. Layout and class: the article's own class is \texttt{article} with packages including amsmath, amsthm, amssymb, amscd; the wrapper is the lane's pinned \texttt{amsart} editorial wrapper at 10pt on A4, which restates the theorem-like environments the retained text needs and adds the article's own macro definitions that the retained text needs (\texttt{\textbackslash forkindep}, \texttt{\textbackslash vd}), with extra wrapper packages (bm, bbm, dsfont, xspace, cancel, mathtools). The delivered numbering is the unit's own. Assets: no retained span contains a figure environment, a graphics-inclusion command, a PDF-page inclusion, a table or a TikZ picture; no image file is copied into this package, the package declares no asset dependency and no image file is redistributed with it. Licence: version arXiv:2606.27011v1 of this article is distributed under the Creative Commons Attribution 4.0 International licence, as recorded in the arXiv licence metadata for this exact version and shown on the abstract page https://arxiv.org/abs/2606.27011v1. A full rights notice is delivered at licenses/arxiv-2606.27011-CC-BY-4.0.txt. Characters: every retained excerpt is pure ASCII, so the delivered engine, XeLaTeX with its default Latin Modern text font, reports no missing character.
\begin{fact}
\label{fact1}
Let $x$ be a set of ordinals.

\begin{enumerate}
\item If $y$ is a set of ordinals and $A\in\vd_x$ is a set of sets of ordinals, then there is $z\in A$ such that $y\forkindep[x]z$;
\item if $A$ is a $\vd_x$-set then $A$ is well-orderable if and only if $A\subset\vd_x$;
\item if $y_0\forkindep[x]y_1$ then $\vd_{xy_0}\cap\vd_{xy_1}=\vd_x$;
\item if $y_0\forkindep[x]y_1$ and $A_0\in\vd_{xy_0}$ and $A_1\in\vd_{xy_1}$ are disjoint subsets of $\vd_x$, then there is a set $B\in\vd_x$ such that $A_0\subseteq B$ and $A_1\cap B=0$.
\end{enumerate}
\end{fact}

\begin{proposition}
\label{ab}
Let $G$ be a Polish group and $C\subset G$ be a nonempty closed set of duplication index $\geq n$. Let $L$ be a distributive Artinian lattice of rank $\leq n$. If $g\in G$ and $c\in C$ are elements mutually Cohen-generic over $x$, then $g\perp^L_xg+c$ holds.
\end{proposition}

\begin{proof}
The argument relies on a general claim of independent interest.

\begin{claim}
\label{duplicationclaim}
Suppose that $\bar y$ is an $n$-tuple $\forkindep$-independent over $x$, $\bar v$ is an $n$-tuple of elements of $L$ such that $\bar v(i)\in \vd_{x\bar y(i)}$, and $u\leq\bigwedge_i\bar v(i)$ is an indecomposable element. Then there is an $i\in n$ and an interpolant between $u$ and $\bar v(i)$ in $\vd_x$.
\end{claim}

\begin{proof}
For every set $b\subseteq n$, let $w_b=\bigwedge\{w\in L\cap \vd_{x\bar y\restriction b}\colon u\leq w\}$. By the Artinian property of $L$, $w_b$ exists as an element of $L$ and belongs to $\vd_{x\bar y\restriction b}$. It must be irreducible: otherwise, the decomposition of it into irreducible elements would belong to $\vd_{x\bar y\restriction b}$, and the irreducibility of $u$ would guarantee that one of the composants of $w_b$ must be above $u$. This would contradict the definition of $w_b$.

Fix $b\subseteq n$. The indepence of $\bar y$ together with Fact~\ref{fact1}(3) shows that the set $\{a\subseteq n\colon w_b\in\vd_{x\bar y\restriction a}\}$ is closed under intersection, so it contains an inclusion-smallest element $a_b$. There are now two cases.

\noindent\textbf{Case 1.} There is a nonempty set $b\subseteq n$ such that $a_b=0$. In such a case, pick $i\in b$ and observe that $w_b$ is an interpolant between $u$ and $\bar v(i)$ in $\vd_x$. This confirms the conclusion of the claim.

\noindent\textbf{Case 2.} Case 1 fails. In this case, by reverse recursion on $i\leq n$ build sets $b_i\subseteq n$ such that $b_n=n$ and $b_{i-1}=b_i$ with any single element of $a_{b_i}$ removed from $b_i$. It is then clear that $w_{b_i}\leq w_{b_{i-1}}$ and the equality does not occur, because $w_{b_i}$ does not belong to $w_{b_{i-1}}$ by the minimal choice of $a_{b_i}$. Thus, $\langle w_{b_i}\colon i\leq n\rangle$ is a strictly decreasing sequence of irreducible elements of $L$ of length $n+1$, contradicting the initial assumptions on $L$. The claim follows.
\end{proof}

Now, let $x$ be a set of ordinals such that $L\in\vd_x$. Let $g\in G$ and $c\in C$ be points mutually Cohen-generic over $x$ in $G$; I must argue that $g\perp^L_x g+c$. In order to do this, let $u_0\in \vd_{xg}$ and $u_1\in\vd_{xg+c}$ be elements of $L$ such that $u_0\leq u_1$; I must find an interpolant in $\vd_x$. In order to do this, let $F_0$, $F_1$ be $\vd$-functions such that $u_0=F_0(x, g)$ and $u_1=F_1(x, g+c)$.

Suppose towards a contradiction that the desired interpolant does not exist. Use the product property of $\forkindep$ to find $\vd_x$ sets $A_0, A_1$ such that $g\in A_0$, $c\in A_1$, and for every $g'\in A_0$ and $c'\in A_1$ such that $g'\forkindep_x c'$, $F_0(x, g')\leq F_1(x, g+c)$ holds and there is no interpolant between them in $\vd_x$.

Let $\bar d\in A_1^n$ be a tuple of elements of $A_1$ which is $\forkindep_{xg}$-independent. By the assumption on the duplication index of $C$, the tuple $\langle g+\bar d(i)\colon i\in n\rangle$ is $\forkindep_x$-independent. By the choice of the set $A_1$, for every $i\in n$, $u_0\leq F_1(x, g+\bar d(i))$ holds and there is no interpolant in $\vd_x$. This, however, immediately contradicts Claim~\ref{duplicationclaim}.
\end{proof}

\end{document}
